\documentclass[12pt]{article}
\usepackage[utf8]{inputenc}
\usepackage[T1]{fontenc}
\usepackage[margin=1in]{geometry}
\usepackage{setspace}
\usepackage{xcolor}
\usepackage[colorlinks=true, urlcolor=blue!60!black, linkcolor=blue!60!black]{hyperref}
\setlength{\parindent}{0em}
\setlength{\parskip}{0.45em}
\renewcommand{\baselinestretch}{1.0}
\begin{document}
\begin{center}
{\Large \textit{Trump v. Anderson}, 601 U.S. 100 (2024)\par}
\vspace{0.3em}
{\normalsize Case Brief --- One-Page Edition\par}
\end{center}

\textbf{Memory Jogger.} The Section~3 ballot case: States cannot disqualify federal candidates (\href{https://www.supremecourt.gov/opinions/23pdf/601us1r06_m6hn.pdf}{601 U.S. 100}).

\textbf{Facts.} Colorado voters led by Norma Anderson sought to keep Donald Trump off the 2024 primary ballot, alleging that he incited the January 6, 2021 breach of the Capitol and was disqualified under Section~3 of the Fourteenth Amendment.

\textbf{Procedural History.} The trial court found insurrection but no coverage of the President; the Colorado Supreme Court, 4--3, ordered Trump off the ballot.

\textbf{Votes.} Reversed, 9--0 in the judgment; a five-Justice per curiam, with Barrett and with Sotomayor, Kagan, and Jackson concurring.

\textbf{Holding.} States may disqualify state officers and candidates but cannot enforce Section~3 against federal offices; five Justices assigned that enforcement to Congress under Section~5.

\textbf{Reasoning.} The Amendment expanded federal power at the States' expense; state exclusions would burden Congress's two-thirds amnesty power; history shows congressional enforcement; and state-by-state rulings would create a presidential patchwork. The concurring Justices objected that the majority decided far more than needed.

\textbf{Examples.} \emph{Same-Side:} a State cannot remove a House or Senate candidate. \emph{Opposite-Side:} a State may bar a former sheriff from its governor's race. \emph{Fence-Sitter:} a State enforcing the two-term limit.

\textbf{Critique.} Requiring legislation lets a bare majority nullify a disability that takes two-thirds of both Houses to lift, and the merits were left unresolved.

\textbf{Key Quotation.} ``We conclude that States may disqualify persons holding or attempting to hold state office. But States have no power under the Constitution to enforce Section 3 with respect to federal offices, especially the Presidency.'' --- Per Curiam, 601 U.S. at 110.

\end{document}
